%% kapa: exam=1 round=35 year=2024 period=1 form=A subject=civil_law subject_ko=민법 seq=5 number=5
%% source: Q-Net 감정평가사 기출문제 (cst00309) artlSeq=5213229 file=raw/1008/5213229/1교시 문제지 원본.pdf
%% conversion: kordoc 4.24.51 -> tools/gen1.py (auto); review=proofread-sonnet; flags=none
\begin{question}{5}
비법인사단 A의 유일한 대표자 甲은 乙에게 대표자로서의 모든 권한을 포괄적으로 위임하고 자신은 이사의 직무를 집행하지 않았다. 이에 관한 설명으로 옳은 것을 모두 고른 것은? (다툼이 있으면 판례에 따름)
\begin{framedtext}
ㄱ. 甲의 행위는 이사의 직무상 선량한 관리자의 주의의무를 위반한 행위이다. \par
ㄴ. 乙이 A의 사실상 대표자로서 丙과 금전소비대차계약을 체결한 경우, 그 계약의 효력은 원칙적으로 A에게 미친다. \par
ㄷ. 乙이 A의 사실상 대표자로서 사무를 집행하면서 그 직무에 관한 불법행위로 丁에게 손해를 입힌 경우, A는 丁에 대하여 법인의 불법행위로 인한 손해배상책임을 부담한다. \par
\end{framedtext}
\begin{choices}
  \item ㄱ
  \item ㄴ
  \item ㄱ, ㄷ
  \item ㄴ, ㄷ
  \item ㄱ, ㄴ, ㄷ
\end{choices}
\end{question}
\answer{3}
